% ============================================================
\setpartpreamble{\partintro
  {Eigene SQL-Ausdrücke, Aggregate, Laufzeit-Metadaten, Scaffolding und ein zweiter
  Provider: die Werkzeuge für sehr komplexe Plugins.}
  {
    \item Kapitel~\ref{ch:custom-sql-expression-nodes}: Eigene SQL-Ausdrucksknoten
    \item Kapitel~\ref{ch:aggregate-translators}: Aggregat-Translators
    \item Kapitel~\ref{ch:property-level-fluent-api-and}: Fluent API auf Property-Ebene und
    Laufzeit-Annotations
    \item Kapitel~\ref{ch:scaffolding-in-practice}: Scaffolding in der Praxis
    \item Kapitel~\ref{ch:what-belongs-in-the-core}: Was in den Kern gehört
    \item Kapitel~\ref{ch:a-second-provider-sql-server}: Ein zweiter Provider: SQL Server
  }}
\part{Tiefe Erweiterungen}\label{part:IX}

% ------------------------------------------------------------
\chapter{Eigene SQL-Ausdrucksknoten}\label{ch:custom-sql-expression-nodes}
\versionrisk{volatile}

Kapitel~\ref{ch:translating-methods-and-members} baut SQL aus vorhandenen Bausteinen:
Funktionen, Konstanten, Casts. Das reicht für die meisten Plugins. Ein \textbf{eigener
Ausdrucksknoten} (eine Ableitung von \code{Sql\allowbreak{}Expression}) wird nötig, wenn
mindestens eins davon zutrifft:

\begin{datatable}{H{0.80}H{1.20}}
  \textbf{Grund} & \textbf{Beispiel} \\
  \tablehead
  Die Syntax ist kein Funktionsaufruf & Operatoren wie \code{@@}, \code{<->},
  \code{IS DISTINCT FROM}, Schlüsselwörter wie \code{AT TIME ZONE} \\
  \tablerule
  Die Nullability-Semantik weicht ab & \code{IS DISTINCT FROM} ist nie NULL, auch wenn ein
  Operand NULL ist \\
  \tablerule
  Andere Teile sollen den Knoten wiedererkennen & Optimierungen, die nur für diesen Ausdruck
  gelten \\
\end{datatable}

Das Beispiel implementiert den null-sicheren Vergleich
\code{EF.\allowbreak{}Functions.\allowbreak{}Is\allowbreak{}Not\allowbreak{}Distinct\allowbreak{}From(\allowbreak{}a,\allowbreak{}
b)} und \code{Is\allowbreak{}Distinct\allowbreak{}From(\allowbreak{}a,\allowbreak{} b)}.
Npgsql bietet ihn nicht als Funktion an; er erfüllt alle drei Gründe.

\section{Anatomie eines Knotens}

\begin{datatable}{H{0.78}H{1.26}H{0.96}}
  \textbf{Mitglied} & \textbf{Zweck} & \textbf{Folge, wenn es fehlt} \\
  \tablehead
  \code{Visit\allowbreak{}Children} / \code{Update} & Visitors können Kinder ersetzen;
  unveränderte Kinder ergeben dieselbe Instanz & EF wirft beim ersten Visitor \\
  \tablerule
  \code{Equals} / \code{Get\allowbreak{}Hash\allowbreak{}Code} & Strukturvergleich für den
  Query-Cache und Optimierungen & Cache-Fehltreffer, falsche Deduplizierung \\
  \tablerule
  \code{Print} & Lesbare Darstellung in DebugView und Fehlermeldungen & Unklare Fehler \\
  \tablerule
  \code{Quote} & C\#-Code für vorkompilierte Abfragen und NativeAOT & Kein AOT-Support \\
  \tablerule
  SQL-Generator & Schreibt den Knoten als SQL & \textbf{Stillschweigend} fehlendes SQL \\
  \tablerule
  Nullability-Prozessor & Kennt die NULL-Semantik & Exception beim Übersetzen \\
\end{datatable}

Die beiden letzten Zeilen betreffen auch die Kombination mit anderen Plugins: Ersetzt ein
zweites Plugin den Generator oder den Prozessor, fehlt das SQL des Knotens still oder die
Abfrage wirft. \code{Composition\allowbreak{}Matrix\allowbreak{}Tests} zeigt beides
(Kapitel~\ref{ch:registering-services}).

% source: src/Acme.EntityFrameworkCore.Tenancy.PostgreSQL/Query/Expressions/DistinctFromExpression.cs
\begin{sourcecode}[Query/\allowbreak{}Expressions/\allowbreak{}Distinct\allowbreak{}From\allowbreak{}Expression.\allowbreak{}cs (\allowbreak{}gekürzt)]{csharp}
public class DistinctFromExpression : SqlExpression
{
    private static ConstructorInfo? _quotingConstructor;

    public DistinctFromExpression(SqlExpression left, SqlExpression right, bool negated,
        RelationalTypeMapping? typeMapping)
        : base(typeof(bool), typeMapping)
    {
        // ... ThrowIfNull ...
        Left = left;
        Right = right;
        IsNegated = negated;
    }

    public virtual SqlExpression Left { get; }
    public virtual SqlExpression Right { get; }
    public virtual bool IsNegated { get; }

    protected override Expression VisitChildren(ExpressionVisitor visitor)
        => Update((SqlExpression)visitor.Visit(Left), (SqlExpression)visitor.Visit(Right));

    public virtual DistinctFromExpression Update(SqlExpression left, SqlExpression right)
        => left == Left && right == Right
            ? this
            : new DistinctFromExpression(left, right, IsNegated, TypeMapping);

    public override Expression Quote()
        => New(
            _quotingConstructor ??= typeof(DistinctFromExpression).GetConstructor(
                [typeof(SqlExpression), typeof(SqlExpression), typeof(bool),
                 typeof(RelationalTypeMapping)])!,
            Left.Quote(), Right.Quote(), Constant(IsNegated),
            RelationalExpressionQuotingUtilities.QuoteTypeMapping(TypeMapping));

    protected override void Print(ExpressionPrinter expressionPrinter)
    {
        expressionPrinter.Visit(Left);
        expressionPrinter.Append(IsNegated ? " IS DISTINCT FROM " : " IS NOT DISTINCT FROM ");
        expressionPrinter.Visit(Right);
    }

    public override bool Equals(object? obj)
        => obj is DistinctFromExpression other
           && (ReferenceEquals(this, other)
               || base.Equals(other) && Left.Equals(other.Left)
                  && Right.Equals(other.Right) && IsNegated == other.IsNegated);

    public override int GetHashCode()
        => HashCode.Combine(base.GetHashCode(), Left, Right, IsNegated);
}
\end{sourcecode}

Der Translator erzeugt den Knoten und gleicht dabei die Type Mappings beider Seiten an: Hat
eine Seite ein Mapping (etwa eine Spalte), bekommt die andere dasselbe. So wird aus einer
Konstante das passende Literal:

% source: src/Acme.EntityFrameworkCore.Tenancy.PostgreSQL/Query/DistinctFromTranslator.cs
\begin{sourcecode}[Query/\allowbreak{}Distinct\allowbreak{}From\allowbreak{}Translator.\allowbreak{}cs (\allowbreak{}Auszug)]{csharp}
var definition = method.GetGenericMethodDefinition();
if (definition != IsNotDistinctFromMethod && definition != IsDistinctFromMethod)
    return null;

var left = arguments[1];                        // arguments[0] ist DbFunctions
var right = arguments[2];
var typeMapping = ExpressionExtensions.InferTypeMapping(left, right)
                  ?? _typeMappingSource.FindMapping(left.Type);

return new DistinctFromExpression(
    _sql.ApplyTypeMapping(left, typeMapping)!,
    _sql.ApplyTypeMapping(right, typeMapping)!,
    negated: definition == IsDistinctFromMethod,
    _typeMappingSource.FindMapping(typeof(bool)));
\end{sourcecode}

\section{Den SQL-Generator erweitern}

Der \code{Query\allowbreak{}Sql\allowbreak{}Generator} schreibt SQL per Visitor. Unbekannte
Knoten gibt er an
\code{Expression\allowbreak{}Visitor.\allowbreak{}Visit\allowbreak{}Extension} weiter, das
nur die Kinder besucht. Ohne eigenen Override entsteht also kein Fehler, sondern \textbf{SQL
ohne den Operator}. Das Plugin erbt vom Npgsql-Generator:

% source: src/Acme.EntityFrameworkCore.Tenancy.PostgreSQL/Query/Internal/TenancyNpgsqlQuerySqlGenerator.cs
\begin{sourcecode}[Query/\allowbreak{}Internal/\allowbreak{}Tenancy\allowbreak{}Npgsql\allowbreak{}Query\allowbreak{}Sql\allowbreak{}Generator.\allowbreak{}cs (\allowbreak{}gekürzt)]{csharp}
public class TenancyNpgsqlQuerySqlGenerator : NpgsqlQuerySqlGenerator
{
    public TenancyNpgsqlQuerySqlGenerator(QuerySqlGeneratorDependencies dependencies,
        IRelationalTypeMappingSource typeMappingSource,
        bool reverseNullOrderingEnabled, Version postgresVersion)
        : base(dependencies, typeMappingSource, reverseNullOrderingEnabled, postgresVersion)
    {
    }

    protected override Expression VisitExtension(Expression extensionExpression)
        => extensionExpression switch
        {
            DistinctFromExpression e => VisitDistinctFrom(e),
            _ => base.VisitExtension(extensionExpression),
        };

    protected virtual Expression VisitDistinctFrom(DistinctFromExpression expression)
    {
        VisitOperand(expression.Left);
        Sql.Append(expression.IsNegated ? " IS DISTINCT FROM " : " IS NOT DISTINCT FROM ");
        VisitOperand(expression.Right);
        return expression;
    }

    // Der Knoten als Operand anderer Operatoren: immer in Klammern
    protected override bool RequiresParentheses(SqlExpression outerExpression,
        SqlExpression innerExpression)
        => innerExpression is DistinctFromExpression
           || base.RequiresParentheses(outerExpression, innerExpression);
}
\end{sourcecode}

Klammern braucht es in zwei Richtungen. \code{Visit\allowbreak{}Operand} klammert
zusammengesetzte Operanden \emph{im} Knoten. Für den Knoten \emph{als} Operand fragt EF
\code{Requires\allowbreak{}Parentheses}, und das kennt nur seine eigenen Knoten. Ohne den
Override schrieb das Plugin
\code{a IS DISTINCT FROM 'x' <> a IS DISTINCT FROM 'y'}; PostgreSQL bindet \code{<>} stärker
als \code{IS} und meldet einen Syntaxfehler. \code{Deep\allowbreak{}Query\allowbreak{}Tests}
prüft das SQL, \code{Advanced\allowbreak{}Tests} die Abfrage gegen PostgreSQL.

Der Generator wird über
\code{IQuery\allowbreak{}Sql\allowbreak{}Generator\allowbreak{}Factory} erzeugt, einen
Singleton-Einzeldienst. Das Plugin ersetzt ihn durch eine eigene Factory mit denselben
Abhängigkeiten wie die von Npgsql
(\code{Query\allowbreak{}Sql\allowbreak{}Generator\allowbreak{}Dependencies},
Type-Mapping-Source, \code{INpgsql\allowbreak{}Singleton\allowbreak{}Options}).

\section{Den Nullability-Prozessor erweitern}

Nach der Übersetzung optimiert EF NULL-Vergleiche. Dafür muss es von jedem Knoten wissen, ob
er NULL werden kann. Unbekannte Knoten landen in
\code{Visit\allowbreak{}Custom\allowbreak{}Sql\allowbreak{}Expression}, dessen
Basisimplementierung eine Exception wirft. Der Prozessor hängt an einer Kette aus drei
Klassen, die alle ersetzt werden müssen:

% source: src/Acme.EntityFrameworkCore.Tenancy.PostgreSQL/Query/Internal/TenancyNpgsqlSqlNullabilityProcessor.cs
\begin{sourcecode}[Query/\allowbreak{}Internal/\allowbreak{}Tenancy\allowbreak{}Npgsql\allowbreak{}Sql\allowbreak{}Nullability\allowbreak{}Processor.\allowbreak{}cs (\allowbreak{}gekürzt)]{csharp}
public class TenancyNpgsqlSqlNullabilityProcessor : NpgsqlSqlNullabilityProcessor
{
    // Konstruktor: (…ProcessorDependencies, …ProcessorParameters) : base(...)

    protected override SqlExpression VisitCustomSqlExpression(
        SqlExpression sqlExpression, bool allowOptimizedExpansion, out bool nullable)
    {
        if (sqlExpression is DistinctFromExpression distinctFrom)
        {
            // Kinder besuchen (sie können optimiert werden), ihre Nullability ist egal
            var left = Visit(distinctFrom.Left, out _)!;
            var right = Visit(distinctFrom.Right, out _)!;

            nullable = false;                   // IS [NOT] DISTINCT FROM ist nie NULL
            return distinctFrom.Update(left, right);
        }

        return base.VisitCustomSqlExpression(
            sqlExpression, allowOptimizedExpansion, out nullable);
    }
}

public class TenancyNpgsqlParameterBasedSqlProcessor : NpgsqlParameterBasedSqlProcessor
{
    protected override Expression ProcessSqlNullability(
        Expression selectExpression, ParametersCacheDecorator parametersDecorator)
        => new TenancyNpgsqlSqlNullabilityProcessor(Dependencies, Parameters)
            .Process(selectExpression, parametersDecorator);
}

// + TenancyNpgsqlParameterBasedSqlProcessorFactory
//   : IRelationalParameterBasedSqlProcessorFactory
\end{sourcecode}

% source: src/Acme.EntityFrameworkCore.Tenancy.PostgreSQL/Infrastructure/TenancyOptionsExtension.cs
\begin{sourcecode}[Infrastructure/\allowbreak{}Tenancy\allowbreak{}Options\allowbreak{}Extension.\allowbreak{}cs › Apply\allowbreak{}Services (\allowbreak{}Ergänzung)]{csharp}
// 4) Tiefe Erweiterung: eigener SQL-Ausdrucksknoten
services.Replace(ServiceDescriptor
    .Singleton<IQuerySqlGeneratorFactory, TenancyNpgsqlQuerySqlGeneratorFactory>());
services.Replace(ServiceDescriptor
    .Scoped<IRelationalParameterBasedSqlProcessorFactory,
        TenancyNpgsqlParameterBasedSqlProcessorFactory>());
\end{sourcecode}

\begin{important}[Der Preis eigener Knoten]
  Eigene Knoten verlangen das Ersetzen zweier Einzeldienste und das Erben von internen
  Npgsql-Klassen. Zwei Plugins, die beide eigene Knoten einführen, lassen sich nur
  kombinieren, wenn eines vom anderen erbt. Deshalb gilt: Funktionskomposition zuerst;
  eigene Knoten nur, wenn einer der drei Gründe vorliegt. Und die Ersetzung gehört an
  prominente Stelle in die Dokumentation.
\end{important}

\begin{note}[Experimentelle API]
  \code{Quote} und
  \code{Relational\allowbreak{}Expression\allowbreak{}Quoting\allowbreak{}Utilities} gehören
  zur experimentellen API für vorkompilierte Abfragen. Ihre Verwendung meldet der Compiler
  als EF9100; das Beispiel unterdrückt die Meldung projektweit und dokumentiert das in
  \code{Directory.\allowbreak{}Build.\allowbreak{}props}.
\end{note}

% source: tests/Acme.EntityFrameworkCore.Tenancy.Tests/DeepQueryTests.cs
\begin{sourcecode}[tests/\allowbreak{}.\allowbreak{}.\allowbreak{}.\allowbreak{}/\allowbreak{}Deep\allowbreak{}Query\allowbreak{}Tests.\allowbreak{}cs (\allowbreak{}Auszug)]{csharp}
[Fact]
public void IsDistinctFrom_in_projection_passes_nullability_processor()
{
    // Ohne eigenen Nullability-Prozessor würde EF hier eine Exception werfen
    using var ctx = Contexts.Create();

    var sql = ctx.Customers
        .Select(c => EF.Functions.IsDistinctFrom(c.Name, (string?)null))
        .ToQueryString();

    Assert.Contains("IS DISTINCT FROM", sql);
}
\end{sourcecode}

\section{Praxisbeispiel: Operatoren für Vektorsuche und pgvector}

Embeddings und Ähnlichkeitssuche bringen Operatoren in SQL, die keine Funktionen sind: pgvector
schreibt die L2-Distanz als \code{<->}, die Kosinus-Distanz als
\code{<=>} und das negative innere Produkt als
\code{<\#>}; die Volltextsuche verwendet \code{@@}.
\code{Sql\allowbreak{}Binary\allowbreak{}Expression} hilft hier nicht, denn ihr Operator ist
ein \code{Expression\allowbreak{}Type}, und für \code{<=>} gibt es
keinen. \code{Custom\allowbreak{}Operator\allowbreak{}Tests} zeigt deshalb
\code{EF.\allowbreak{}Functions.\allowbreak{}Cosine\allowbreak{}Distance(\allowbreak{}a,\allowbreak{}
b)} als eigenen Knoten \code{Vector\allowbreak{}Distance\allowbreak{}Expression} mit dem
Operator als Zeichenkette, mit Translator, Generator und Nullability-Prozessor. Es ist
Testcode: Eine Anwendung nimmt \code{Pgvector.\allowbreak{}Entity\allowbreak{}Framework\allowbreak{}Core},
das den Typ \code{vector}, sein Type Mapping und diese Übersetzungen mitbringt. Npgsql selbst
hat für solche Fälle \code{Pg\allowbreak{}Unknown\allowbreak{}Binary\allowbreak{}Expression},
einen Knoten mit beliebigem Operator und propagierender Nullability, der Generator und
Prozessor nicht ersetzen muss. Er liegt aber in einem \code{Internal}-Namespace (EF1001).

Drei Dinge entscheiden darüber, ob ein solcher Operator richtiges SQL ergibt.

\textbf{Type Mappings.} Der Translator gibt beiden Operanden dasselbe Mapping, abgeleitet von
der Seite, die eins hat, meist der Spalte. So geht der Parameter mit dem Speichertyp der Spalte
an die Datenbank:

% source: tests/Acme.EntityFrameworkCore.Tenancy.Tests/CustomOperatorTests.cs
\begin{sourcecode}[tests/\allowbreak{}.\allowbreak{}.\allowbreak{}.\allowbreak{}/\allowbreak{}Custom\allowbreak{}Operator\allowbreak{}Tests.\allowbreak{}cs ›\allowbreak{} Translator]{csharp}
if (method != CosineDistanceMethod)
    return null;

var left = arguments[1];                        // arguments[0] ist DbFunctions
var right = arguments[2];
var operandMapping = ExpressionExtensions.InferTypeMapping(left, right)
                     ?? _typeMappingSource.FindMapping(left.Type);

return new VectorDistanceExpression(
    _sql.ApplyTypeMapping(left, operandMapping)!,
    _sql.ApplyTypeMapping(right, operandMapping)!,
    "<=>",
    _typeMappingSource.FindMapping(typeof(double)));
\end{sourcecode}

Im Test ist die Spalte ein \code{float[]}, also \code{real[]}. pgvector definiert seine
Operatoren aber nur für \code{vector}; gegen eine echte Datenbank braucht es deshalb ein Type
Mapping mit diesem Speichertyp, wie es das Paket registriert. Das Ergebnis bekommt das Mapping
von \code{double}, also \code{double precision}.

\textbf{Rangfolge und Klammern.} EF kennt die Rangfolge eines fremden Operators nicht.
PostgreSQL ordnet selbst definierte Operatoren zwischen die arithmetischen und die
Vergleichsoperatoren ein: \code{a <=> b * 2} heißt
\code{a <=> (\allowbreak{}b * 2)}. Ohne Gegenmaßnahme schrieb der
Generator für \code{Cosine\allowbreak{}Distance(\allowbreak{}.\allowbreak{}.\allowbreak{}.\allowbreak{})
* 2 + 1} genau das. Die Lösung sind dieselben zwei Richtungen wie oben: Operanden im Knoten
klammern und den Knoten als Operand über
\code{Requires\allowbreak{}Parentheses} immer klammern. \code{Print} klammert genauso, damit
die DebugView zeigt, was die Datenbank liest.

% source: tests/Acme.EntityFrameworkCore.Tenancy.Tests/CustomOperatorTests.cs
\begin{sourcecode}[tests/\allowbreak{}.\allowbreak{}.\allowbreak{}.\allowbreak{}/\allowbreak{}Custom\allowbreak{}Operator\allowbreak{}Tests.\allowbreak{}cs ›\allowbreak{} Generator]{csharp}
public class VectorQuerySqlGenerator : TenancyNpgsqlQuerySqlGenerator
{
    // ... Konstruktor, VisitExtension, VisitOperand

    protected virtual Expression VisitVectorDistance(VectorDistanceExpression expression)
    {
        VisitOperand(expression.Left);
        Sql.Append($" {expression.Operator} ");
        VisitOperand(expression.Right);
        return expression;
    }

    // Ohne Override: "a <=> b * 2.0", gelesen als a <=> (b * 2.0)
    protected override bool RequiresParentheses(SqlExpression outerExpression,
        SqlExpression innerExpression)
        => innerExpression is VectorDistanceExpression
           || base.RequiresParentheses(outerExpression, innerExpression);
}
\end{sourcecode}

Der Generator erbt vom Generator des Plugins, nicht von dem von Npgsql: Beide ersetzen
denselben Einzeldienst, und nur so bleiben beide Knoten übersetzbar. Das ist der Preis aus
dem Kasten oben, im Kleinen.

\textbf{Nullability.} Ein gewöhnlicher Operator ist NULL, sobald ein Operand NULL ist. Der
Prozessor meldet das weiter:

% source: tests/Acme.EntityFrameworkCore.Tenancy.Tests/CustomOperatorTests.cs
\begin{sourcecode}[tests/\allowbreak{}.\allowbreak{}.\allowbreak{}.\allowbreak{}/\allowbreak{}Custom\allowbreak{}Operator\allowbreak{}Tests.\allowbreak{}cs ›\allowbreak{} Nullability]{csharp}
if (sqlExpression is VectorDistanceExpression distance)
{
    // NULL hinein, NULL heraus
    var left = Visit(distance.Left, out var leftNullable)!;
    var right = Visit(distance.Right, out var rightNullable)!;

    nullable = leftNullable || rightNullable;
    return distance.Update(left, right);
}
\end{sourcecode}

Die Meldung hat sichtbare Folgen. Für
\code{!(\allowbreak{}Cosine\allowbreak{}Distance(\allowbreak{}.\allowbreak{}.\allowbreak{}.\allowbreak{})
< 0.5)} schreibt EF bei einer Pflichtspalte \code{>= 0.5}, bei einer optionalen Spalte ein
\code{CASE}, weil \code{NOT (\allowbreak{}NULL < 0.5)} in SQL NULL ist, in C\# aber
\code{true}. Ohne den Prozessor scheitert die Übersetzung. Die Tests prüfen das SQL aller
Fälle:

% source: tests/Acme.EntityFrameworkCore.Tenancy.Tests/CustomOperatorTests.cs
\begin{sourcecode}[tests/\allowbreak{}.\allowbreak{}.\allowbreak{}.\allowbreak{}/\allowbreak{}Custom\allowbreak{}Operator\allowbreak{}Tests.\allowbreak{}cs (\allowbreak{}Auszug)]{csharp}
var sql = ctx.Documents
    .Select(d => EF.Functions.CosineDistance(d.Embedding, query) * 2 + 1)
    .ToQueryString();

Assert.Contains("SELECT (d.\"Embedding\" <=> @query) * 2.0 + 1.0", sql);
\end{sourcecode}

% ------------------------------------------------------------
\chapter{Aggregat-Translators}\label{ch:aggregate-translators}
\versionrisk{medium}

Aggregate wie \code{Sum} oder \code{string\_\allowbreak{}agg} fassen mehrere Zeilen
zusammen, typischerweise nach \code{Group\allowbreak{}By}. Ein
\code{IAggregate\allowbreak{}Method\allowbreak{}Call\allowbreak{}Translator} bekommt statt
eines einzelnen Arguments eine \code{Enumerable\allowbreak{}Expression}, die beschreibt,
\emph{was} aggregiert wird:

\begin{datatable}{H{0.70}H{1.09}H{1.21}}
  \textbf{Eigenschaft} & \textbf{Bedeutung} & \textbf{Herkunft in LINQ} \\
  \tablehead
  \code{Selector} & Der aggregierte Ausdruck & \code{g.\allowbreak{}Select(\allowbreak{}o =>
  o.\allowbreak{}Total > 0)} \\
  \tablerule
  \code{Is\allowbreak{}Distinct} & Nur verschiedene Werte &
  \code{g.\allowbreak{}Select(\allowbreak{}.\allowbreak{}.\allowbreak{}.\allowbreak{}).\allowbreak{}Distinct(\allowbreak{})}
  \\
  \tablerule
  \code{Predicate} & Filter vor dem Aggregat & \code{g.\allowbreak{}Where(\allowbreak{}o =>
  o.\allowbreak{}Id > 10)} \\
  \tablerule
  \code{Orderings} & Reihenfolge (für reihenfolgeabhängige Aggregate) &
  \code{g.\allowbreak{}Order\allowbreak{}By(\allowbreak{}.\allowbreak{}.\allowbreak{}.\allowbreak{})}
  \\
\end{datatable}

Das Beispiel ergänzt
\code{EF.\allowbreak{}Functions.\allowbreak{}Every(\allowbreak{}values)} für PostgreSQLs
\code{bool\_\allowbreak{}and}:

% source: src/Acme.EntityFrameworkCore.Tenancy.PostgreSQL/Query/EveryAggregateTranslator.cs
\begin{sourcecode}[Query/\allowbreak{}Every\allowbreak{}Aggregate\allowbreak{}Translator.\allowbreak{}cs (\allowbreak{}Auszug)]{csharp}
public virtual SqlExpression? Translate(MethodInfo method, EnumerableExpression source,
    IReadOnlyList<SqlExpression> arguments,
    IDiagnosticsLogger<DbLoggerCategory.Query> logger)
{
    if (method != EveryMethod || source.Selector is not SqlExpression selector)
        return null;

    // Npgsqls Factory setzt DISTINCT, FILTER (WHERE ...) und ORDER BY im Aggregat um
    if (_sql is not NpgsqlSqlExpressionFactory npgsql)
        return null;

    var boolean = _typeMappingSource.FindMapping(typeof(bool))!;

    // bool_and liefert für eine leere Menge NULL, Enumerable.All liefert true
    return _sql.Coalesce(
        npgsql.AggregateFunction("bool_and", [selector], source,
            nullable: true, argumentsPropagateNullability: [false],
            typeof(bool), boolean),
        _sql.Constant(true, boolean),
        boolean);
}
\end{sourcecode}

\begin{sourcecode}[Erzeugtes SQL für g.\allowbreak{}Where(\allowbreak{}o => o.\allowbreak{}Id > 10).\allowbreak{}Select(\allowbreak{}o => o.\allowbreak{}Total > 0)]{sql}
SELECT COALESCE(bool_and(o."Total" > 0.0) FILTER (WHERE o."Id" > 10), TRUE)
FROM "Orders" AS o
GROUP BY o."TenantId"
\end{sourcecode}

\begin{important}[Semantik vor Syntax]
  Ein Aggregat-Translator muss die Semantik der .NET-Seite treffen, nicht nur die Syntax der
  Datenbank. Typische Fallen: leere Mengen (NULL statt 0, true oder Leerstring), NULL-Werte
  in der Eingabe und Prädikate. Npgsql selbst gleicht etwa \code{string.\allowbreak{}Join}
  mit \code{COALESCE} an \code{string\_\allowbreak{}agg} an; das Beispiel folgt diesem
  Muster.
\end{important}

Datenbanken ohne \code{FILTER}-Klausel (etwa SQL Server) emulieren Prädikate mit \code{CASE
WHEN prädikat THEN wert END}. Ein provider-neutraler Aggregat-Translator muss beides kennen;
deshalb liegt dieser im PostgreSQL-Paket.

% ------------------------------------------------------------
\chapter{Fluent API auf Property-Ebene und Laufzeit-Annotations}\label{ch:property-level-fluent-api-and}
\versionrisk{volatile}

\section{Property-Builder}

Bisher setzt das Plugin voraus, dass die Mandanten-Property \code{Tenant\allowbreak{}Id}
heißt und \code{ITenant\allowbreak{}Owned} implementiert ist. Mit einer Fluent-API-Methode
auf \code{Property\allowbreak{}Builder} geht es ohne beides, etwa bei bestehenden Modellen
mit einer Property \code{Mandant\allowbreak{}Id} oder einer expliziten
Interface-Implementierung:

\begin{sourcecode}[Anwendung]{csharp}
modelBuilder.Entity<Contract>().Property(c => c.Mandant).IsTenantProperty();
\end{sourcecode}

% source: src/Acme.EntityFrameworkCore.Tenancy/Metadata/TenancyPropertyBuilderExtensions.cs
\begin{sourcecode}[Metadata/\allowbreak{}Tenancy\allowbreak{}Property\allowbreak{}Builder\allowbreak{}Extensions.\allowbreak{}cs (\allowbreak{}Kernpaket)]{csharp}
public static PropertyBuilder IsTenantProperty(
    this PropertyBuilder builder, bool isTenant = true)
{
    ArgumentNullException.ThrowIfNull(builder);

    if (isTenant)
        builder.Metadata.SetAnnotation(TenancyAnnotationNames.IsTenantProperty, true);
    else
        builder.Metadata.RemoveAnnotation(TenancyAnnotationNames.IsTenantProperty);

    return builder;
}

public static PropertyBuilder<TProperty> IsTenantProperty<TProperty>(
    this PropertyBuilder<TProperty> builder, bool isTenant = true)
    => (PropertyBuilder<TProperty>)IsTenantProperty((PropertyBuilder)builder, isTenant);

public static IConventionPropertyBuilder? IsTenantProperty(
    this IConventionPropertyBuilder builder,
    bool isTenant = true,
    bool fromDataAnnotation = false)
    => builder.HasAnnotation(TenancyAnnotationNames.IsTenantProperty,
           isTenant ? true : null, fromDataAnnotation) is null
        ? null
        : builder;
\end{sourcecode}

Die Methode liegt im Kernpaket, weil sie nur Metadaten setzt. Die Conventions beider
Provider-Pakete suchen die Mandanten-Property über dieselbe Kernfunktion:

\begin{sourcecode}[Metadata/\allowbreak{}Tenancy\allowbreak{}Entity\allowbreak{}Type\allowbreak{}Extensions.\allowbreak{}cs (\allowbreak{}Kernpaket,\allowbreak{} gekürzt, vereinfacht)]{csharp}
public static IReadOnlyProperty? FindTenantPropertyCandidate(
    this IReadOnlyEntityType entityType)
{
    var marked = entityType.GetProperties()
        .Where(p => p[TenancyAnnotationNames.IsTenantProperty] is true)
        .ToList();

    if (marked.Count > 1)
        throw new InvalidOperationException("... mehrere Mandanten-Properties ...");
    if (marked.Count == 1)
        return marked[0];

    return typeof(ITenantOwned).IsAssignableFrom(entityType.ClrType)
        ? entityType.FindProperty(nameof(ITenantOwned.TenantId))
        : null;
}
\end{sourcecode}

\begin{datatable}{H{0.92}H{1.10}H{0.98}}
  \textbf{Builder} & \textbf{Typische Plugin-Methoden} & \textbf{Hinweis} \\
  \tablehead
  \code{Property\allowbreak{}Builder} / \code{<\allowbreak{}T>} & Spaltenoptionen, Rollen
  von Properties & Hier: \code{Is\allowbreak{}Tenant\allowbreak{}Property} \\
  \tablerule
  \code{Index\allowbreak{}Builder} / \code{<\allowbreak{}T>} & Index-Methoden,
  Operator-Klassen, Filter & Annotation an \code{Metadata} (IMutableIndex) \\
  \tablerule
  \code{Key\allowbreak{}Builder} & Schlüsseloptionen (Clustering) & Wie oben \\
  \tablerule
  \code{Complex\allowbreak{}Type\allowbreak{}Property\allowbreak{}Builder} & Properties in
  Complex Types & Nur anbieten, wenn die Option dort Sinn ergibt
  (Kapitel~\ref{ch:complex-types-and-json}) \\
  \tablerule
  \code{Model\allowbreak{}Builder} & Modellweite Optionen & Annotation am Modell \\
\end{datatable}

\section{Laufzeit-Annotations}

Der SaveChanges-Interceptor muss zur Laufzeit wissen, welche Entitäten isoliert sind und wie
ihre Mandanten-Property heißt. Die Convention schreibt das in eine Annotation am Entity-Typ;
der Interceptor liest sie aus dem \emph{Laufzeitmodell}:

\begin{sourcecode}[Convention und Kernpaket (\allowbreak{}Auszug)]{csharp}
// Convention (Design-Time-Modell)
entityType.Builder.HasAnnotation(TenancyAnnotationNames.TenantProperty, property.Name);

// Interceptor (Laufzeitmodell)
public static string? GetTenantPropertyName(this IReadOnlyEntityType entityType)
    => entityType.GetRootType()[TenancyAnnotationNames.TenantProperty] as string;
\end{sourcecode}

Zwei Fragen entscheiden, ob das funktioniert. Landet die Annotation im Laufzeitmodell? Ja:
\code{Runtime\allowbreak{}Model\allowbreak{}Convention} entfernt nur EFs eigene
Design-Time-Annotations und übernimmt alle anderen. Landet sie im Compiled Model? Ja:
Einfache Werte schreibt EF automatisch
(Kapitel~\ref{ch:compiled-models-trimming-and-nativeaot}). Der eigene
\code{ICSharp\allowbreak{}Runtime\allowbreak{}Annotation\allowbreak{}Code\allowbreak{}Generator}
hat deshalb die umgekehrte Aufgabe: Er entfernt die RLS-Annotations, die nur Migrations
brauchen:

\begin{sourcecode}[Design/\allowbreak{}Internal/\allowbreak{}Tenancy\allowbreak{}CSharp\allowbreak{}Runtime\allowbreak{}Annotation\allowbreak{}Code\allowbreak{}Generator.\allowbreak{}cs (\allowbreak{}gekürzt, vereinfacht)]{csharp}
public class TenancyCSharpRuntimeAnnotationCodeGenerator
    : NpgsqlCSharpRuntimeAnnotationCodeGenerator
{
    // Konstruktor: (CSharpRuntimeAnnotationCodeGeneratorDependencies,
    //               RelationalCSharpRuntimeAnnotationCodeGeneratorDependencies) : base(...)

    public override void Generate(IEntityType entityType,
        CSharpRuntimeAnnotationCodeGeneratorParameters parameters)
    {
        if (!parameters.IsRuntime)
        {
            foreach (var key in parameters.Annotations.Keys
                         .Where(k => k.StartsWith(RlsAnnotationNames.Prefix,
                                                  StringComparison.Ordinal))
                         .ToList())
            {
                parameters.Annotations.Remove(key);
            }
        }

        base.Generate(entityType, parameters);
    }
}
\end{sourcecode}

Wird das Compiled Model für NativeAOT erzeugt
(\code{For\allowbreak{}Native\allowbreak{}Aot}), schreibt EF zusätzlich das relationale
Modell. Der Annotation-Provider hat die RLS-Annotations auf die Tabellen kopiert, sie
tauchen dort also ein zweites Mal auf. Der Generator wendet deshalb denselben Filter in
einer Überschreibung von \code{Generate(\allowbreak{}ITable,\allowbreak{}
.\allowbreak{}.\allowbreak{}.\allowbreak{})} an.

\begin{datatable}{H{1.36}H{0.88}H{0.88}H{0.88}}
  \textbf{Annotation} & \textbf{Design-Time} & \textbf{Laufzeit} & \textbf{Compiled Model} \\
  \tablehead
  \code{Acme:Rls:*} (Policies) & ja (Migrations) & ja, ungenutzt & gefiltert \\
  \tablerule
  \code{Acme:Tenancy:Tenant\allowbreak{}Property} & ja & ja (Interceptor) & ja \\
  \tablerule
  \code{Acme:Tenancy:Is\allowbreak{}Tenant\allowbreak{}Property} & ja (Convention) & ja,
  ungenutzt & ja (klein) \\
\end{datatable}

\begin{note}[Laufzeit-Annotation und Runtime-Annotation]
  „Laufzeit-Annotation“ heißt in diesem Buch: eine gewöhnliche Annotation, die zur Laufzeit
  gelesen wird. EF kennt zusätzlich eine eigene Runtime-Annotation-API
  (\code{Add\allowbreak{}Runtime\allowbreak{}Annotation},
  \code{Get\allowbreak{}Or\allowbreak{}Add\allowbreak{}Runtime\allowbreak{}Annotation\allowbreak{}Value}
  auf \code{IAnnotatable}). Sie legt Werte erst am fertigen Modell ab, etwa abgeleitete
  Daten, die nicht bei jedem Zugriff neu berechnet werden sollen. Solche Werte erscheinen nie
  in Snapshot oder Migrations. Für Konfiguration, die ein Plugin beim Modellbau festlegt, ist
  die gewöhnliche Annotation der richtige Weg.
\end{note}

\begin{tip}
  Für jede Annotation sollte klar sein, wer sie wann liest. Nur zur Design-Zeit gelesene
  Annotations gehören nicht ins Compiled Model; zur Laufzeit gelesene müssen hinein und
  brauchen einen Test gegen das generierte Modell. Der Compiled-Model-Test des Beispiels
  prüft beides.
\end{tip}

% ------------------------------------------------------------
\chapter{Scaffolding in der Praxis}\label{ch:scaffolding-in-practice}
\versionrisk{volatile}

Beim Scaffolding (\code{dotnet ef dbcontext scaffold}) liest EF das Schema und erzeugt
Entitäten und Context. Ein Plugin, das eigene Schema-Objekte verwaltet, sollte sie dabei
erkennen. Der Weg führt über drei Stationen, von denen das Plugin zwei bereits kennt:

\begin{flowdiagram}
  \flowstep{IDatabaseModelFactory}{Liest das Schema in ein DatabaseModel; das Plugin hängt
  RLS-Daten als Annotations an die Tabellen}
  \flowstep{Scaffolding-Model-Factory}{EF überträgt Tabellen-Annotations auf die
  Entity-Typen}
  \flowstep{IAnnotationCodeGenerator}{Schreibt sie als HasRowLevelSecurity/HasPolicy
  (Kapitel~\ref{ch:design-time-services})}
\end{flowdiagram}

% source: src/Acme.EntityFrameworkCore.Tenancy.PostgreSQL/Scaffolding/TenancyDatabaseModelFactory.cs
\begin{sourcecode}[Scaffolding/\allowbreak{}Tenancy\allowbreak{}Database\allowbreak{}Model\allowbreak{}Factory.\allowbreak{}cs (\allowbreak{}gekürzt)]{csharp}
public class TenancyDatabaseModelFactory : NpgsqlDatabaseModelFactory
{
    public TenancyDatabaseModelFactory(
        IDiagnosticsLogger<DbLoggerCategory.Scaffolding> logger)
        : base(logger)
    {
    }

    public override DatabaseModel Create(
        DbConnection dbConnection, DatabaseModelFactoryOptions options)
    {
        var model = base.Create(dbConnection, options);
        // ... Verbindung öffnen, falls nötig ...
        ReadRowLevelSecurity(dbConnection, model);    // pg_class.relrowsecurity / ...force...
        ReadPolicies(dbConnection, model);            // pg_policies
        return model;
    }

    protected virtual void ReadPolicies(DbConnection connection, DatabaseModel model)
    {
        // SELECT schemaname, tablename, policyname, permissive, roles::text[], cmd,
        //        qual, with_check FROM pg_catalog.pg_policies
        // ...
        var policy = new RlsPolicy(
            Name: reader.GetString(2),
            Command: Enum.Parse<PolicyCommand>(reader.GetString(5), ignoreCase: true),
            Permissive: reader.GetString(3) == "PERMISSIVE",
            Roles: roles is ["public"] ? null : roles,    // "public" = Default ohne TO
            Using: reader.IsDBNull(6) ? null : reader.GetString(6),
            WithCheck: reader.IsDBNull(7) ? null : reader.GetString(7));

        table[RlsAnnotationNames.PolicyPrefix + policy.Name] = policy.Serialize();
    }
}
\end{sourcecode}

Registriert wird die Factory wie der Code-Generator in den Design-Time-Services mit
\code{Remove\allowbreak{}All}, aber mit \code{Add\allowbreak{}Scoped}: EF und Npgsql
registrieren die Database-Model-Factory als Scoped, und sie hängt vom ebenfalls Scoped
Scaffolding-Logger ab. Das Datenformat ist dasselbe wie bei der Fluent API; deshalb funktionieren Snapshot, Migrations und
Code-Generator ohne weitere Anpassung.

\begin{note}[Normalisierte Ausdrücke]
  PostgreSQL speichert Policy-Ausdrücke in normalisierter Form, etwa
  \code{(\allowbreak{}(\allowbreak{}\textquotedbl{}Tenant\allowbreak{}Id\textquotedbl{} =
  (\allowbreak{}NULLIF(\allowbreak{}current\_\allowbreak{}setting(\allowbreak{}'app.\allowbreak{}tenant\_\allowbreak{}id'::text,\allowbreak{}
  true),\allowbreak{} ''::text))::uuid))}. Gescaffoldete Policies enthalten genau diesen
  Text. Das ist korrekt, sieht aber anders aus als die von der Convention erzeugte Form.
  Teams, die nach dem Scaffolding auf Migrations umsteigen, sollten die Policies einmal
  vereinheitlichen, sonst erzeugt die erste Migration ein DROP/CREATE mit gleichem Inhalt.
\end{note}

% source: tests/Acme.EntityFrameworkCore.Tenancy.IntegrationTests/AdvancedTests.cs
\begin{sourcecode}[tests/\allowbreak{}.\allowbreak{}.\allowbreak{}.\allowbreak{}Integration\allowbreak{}Tests/\allowbreak{}Advanced\allowbreak{}Tests.\allowbreak{}cs (\allowbreak{}Auszug)]{csharp}
[Fact]
public void Scaffolding_reads_rls_and_policies()
{
    using var ctx = ContextFor(null);
    var factory = new TenancyDatabaseModelFactory(
        ctx.GetService<IDiagnosticsLogger<DbLoggerCategory.Scaffolding>>());

    var model = factory.Create(
        _fixture.OwnerConnectionString, new DatabaseModelFactoryOptions());
    var orders = model.Tables.Single(t => t.Name == "Orders");

    Assert.Equal(true, orders[RlsAnnotationNames.Enabled]);
    Assert.Equal(true, orders[RlsAnnotationNames.Forced]);
}
\end{sourcecode}

% ------------------------------------------------------------
\chapter{Was in den Kern gehört}\label{ch:what-belongs-in-the-core}
\versionrisk{stable}

Ein zweiter Provider ist der Härtetest für die Paketarchitektur. Er zeigt, welche Teile
wirklich provider-unabhängig sind. Die Faustregel: \textbf{Was keine SQL-Syntax und keinen
Provider-Typ kennt, gehört in den Kern.} Für das Beispiel ergibt sich:

\begin{datatable}{H{1.21}H{0.45}H{1.34}}
  \textbf{Baustein} & \textbf{Ort} & \textbf{Begründung} \\
  \tablehead
  \code{Tenant\allowbreak{}Id}, Konverter, JSON-Reader/Writer & Kern & Reine .NET-Typen \\
  \tablerule
  \code{ITenant\allowbreak{}Owned}, \code{ITenant\allowbreak{}Provider}, Attribute & Kern &
  Für Entitäts-Assemblies ohne Provider-Abhängigkeit \\
  \tablerule
  \code{Is\allowbreak{}Tenant\allowbreak{}Property(\allowbreak{})}, Kandidatensuche & Kern &
  Nur Metadaten \\
  \tablerule
  Laufzeit-Annotation, \code{Get\allowbreak{}Tenant\allowbreak{}Property\allowbreak{}Name} &
  Kern & Vom provider-neutralen Interceptor gelesen \\
  \tablerule
  SaveChanges-Interceptor & Kern & Arbeitet nur mit dem Change Tracker \\
  \tablerule
  \code{Tenant\allowbreak{}Session\allowbreak{}Interceptor} (abstrakt) & Kern & Ablauf
  gleich, nur der SQL-Befehl unterscheidet sich \\
  \tablerule
  Evaluatable-Filter mit Attribut & Kern & Erkennt Funktionsklassen aller Pakete \\
  \tablerule
  Prüfung beim Laden (Materialization-Interceptor) & Kern & Liest nur Modell, Context und
  Werte der Zeile \\
  \tablerule
  Value Generation, Logging, \code{Add\allowbreak{}Acme\allowbreak{}Tenancy} & Kern & Ohne
  SQL-Bezug \\
  \tablerule
  \code{Database.\allowbreak{}Apply\allowbreak{}Tenant\allowbreak{}Session(\allowbreak{})} &
  Kern & Nutzt den Interceptor des Providers \\
  \tablerule
  Type Mapping, Translators, Migrations, RLS-Conventions & Provider & Store-Typen,
  SQL-Syntax, Schema-Objekte \\
  \tablerule
  Eigene SQL-Knoten, Aggregate & Provider & Abhängig von Syntax und Provider-Factory \\
\end{datatable}

\begin{important}[Namenskonflikte zwischen Provider-Paketen]
  Zwei Provider-Pakete, die beide eine Extension-Method
  \code{Use\allowbreak{}Acme\allowbreak{}Tenancy} auf
  \code{Db\allowbreak{}Context\allowbreak{}Options\allowbreak{}Builder} definieren, sind
  nicht gleichzeitig referenzierbar: Der Aufruf wäre mehrdeutig. Das PostgreSQL-Paket nutzt
  den Builder von EF, das SQL-Server-Paket den Builder des Providers
  (\code{Use\allowbreak{}Sql\allowbreak{}Server(\allowbreak{}cs,\allowbreak{} x =>
  x.\allowbreak{}Use\allowbreak{}Acme\allowbreak{}Tenancy(\allowbreak{}))}). Wer mehrere
  Provider-Pakete plant, sollte von Anfang an konsequent den Provider-Builder verwenden, wie
  das NetTopologySuite-Plugin.
\end{important}

% ------------------------------------------------------------
\chapter{Ein zweiter Provider: SQL Server}\label{ch:a-second-provider-sql-server}
\versionrisk{medium}

Das Paket
\code{Acme.\allowbreak{}Entity\allowbreak{}Framework\allowbreak{}Core.\allowbreak{}Tenancy.\allowbreak{}Sql\allowbreak{}Server}
setzt dieselbe Idee mit den Mitteln von SQL Server um. Die Unterschiede zeigen, was ein
Plugin pro Datenbank tatsächlich neu lösen muss:

\begin{datatable}{H{0.81}H{1.00}H{1.19}}
  \textbf{Aspekt} & \textbf{PostgreSQL} & \textbf{SQL Server} \\
  \tablehead
  Mandant in der Session & \code{set\_\allowbreak{}config(\allowbreak{}name,\allowbreak{}
  wert,\allowbreak{} false)} &
  \code{sp\_\allowbreak{}set\_\allowbreak{}session\_\allowbreak{}context} \\
  \tablerule
  \code{Current\allowbreak{}Tenant(\allowbreak{})} in SQL &
  \code{nullif(\allowbreak{}current\_\allowbreak{}setting(\allowbreak{}.\allowbreak{}.\allowbreak{}.\allowbreak{}),\allowbreak{}
  '')::uuid} & \code{CAST(\allowbreak{}SESSION\_\allowbreak{}CONTEXT(\allowbreak{}N'…') AS
  uniqueidentifier)} \\
  \tablerule
  Spaltentyp & \code{uuid} & \code{uniqueidentifier} \\
  \tablerule
  RLS-Mechanismus & \code{CREATE POLICY} pro Tabelle & Predicate-Funktion + \code{CREATE
  SECURITY POLICY} \\
  \tablerule
  Lesen / Schreiben & \code{USING} / \code{WITH CHECK} & \code{FILTER} / \code{BLOCK
  .\allowbreak{}.\allowbreak{}.\allowbreak{} AFTER INSERT/\allowbreak{}UPDATE} \\
  \tablerule
  Tabelleneigentümer & Nur mit \code{FORCE} gebunden & Kein FORCE; Prädikate sind Teil der
  Abfrage \\
  \tablerule
  \code{DROP TABLE}, Umbenennen & Policies verschwinden mit bzw. folgen & Policy muss vorher weg und danach neu angelegt werden (SCHEMABINDING) \\
  \tablerule
  DDL in Batches & Transaktional, beliebig kombinierbar & \code{CREATE
  FUNCTION/\allowbreak{}SCHEMA} allein im Batch:
  \code{EXEC(\allowbreak{}N'.\allowbreak{}.\allowbreak{}.\allowbreak{}')} \\
  \tablerule
  Tabellenabbildungen & \multicolumn{2}{>{\RaggedRight\arraybackslash}p{0.6\linewidth}}{Gleicher
  Vertrag aus dem Kernpaket (Kapitel~\ref{ch:conventions}, Seite~\pageref{sec:supported-table-mappings})} \\
  \tablerule
  Bewusste Ausnahme & Attribut oder
  \code{Has\allowbreak{}No\allowbreak{}Row\allowbreak{}Level\allowbreak{}Security(\allowbreak{})}
  & Nur das Attribut \\
\end{datatable}

\section{Aktivierung über den Provider-Builder}

\begin{sourcecode}[Anwendung und Infrastructure/\allowbreak{}Sql\allowbreak{}Server\allowbreak{}Tenancy\allowbreak{}Options\allowbreak{}Builder\allowbreak{}Extensions.\allowbreak{}cs (\allowbreak{}gekürzt, vereinfacht)]{csharp}
services.AddAcmeTenancy<HttpTenantProvider>();                    // Kernpaket
services.AddDbContext<ShopContext>((sp, o) => o
    .UseSqlServer(connectionString, sql => sql.UseAcmeTenancy(sp)));

public static SqlServerDbContextOptionsBuilder UseAcmeTenancy(
    this SqlServerDbContextOptionsBuilder sqlServerOptionsBuilder,
    string sessionKey = "TenantId",
    string policySchema = "acme")
{
    // Vom Provider-Builder zum eigentlichen DbContextOptionsBuilder
    var optionsBuilder =
        ((IRelationalDbContextOptionsBuilderInfrastructure)sqlServerOptionsBuilder)
        .OptionsBuilder;

    var extension = optionsBuilder.Options.FindExtension<SqlServerTenancyOptionsExtension>()
                    ?? new SqlServerTenancyOptionsExtension();
    // ... Settings validieren und setzen ...
    ((IDbContextOptionsBuilderInfrastructure)optionsBuilder)
        .AddOrUpdateExtension(extension.WithSettings(settings));

    return sqlServerOptionsBuilder;
}
\end{sourcecode}

\section{Migrations: Security Policy statt Policy-Ausdruck}

Die Convention vermerkt für SQL Server nur die Mandantenspalte; das SQL entsteht im
Generator. Die Predicate-Funktion ist für alle Tabellen dieselbe und wird idempotent
angelegt:

\begin{sourcecode}[Migrations/\allowbreak{}Internal/\allowbreak{}Sql\allowbreak{}Server\allowbreak{}Tenancy\allowbreak{}Migrations\allowbreak{}Sql\allowbreak{}Generator.\allowbreak{}cs (\allowbreak{}Auszug, vereinfacht)]{csharp}
protected virtual void CreatePolicy(string table, string? schema, string column,
    MigrationCommandListBuilder builder)
{
    var function = Delimit(PredicateFunctionName, Settings.PolicySchema);
    var target = Delimit(table, schema);
    var predicate = $"{function}({Delimit(column)})";

    // CREATE SCHEMA/FUNCTION müssen allein im Batch stehen -> EXEC(N'...')
    Statement(builder, $"IF SCHEMA_ID({Literal(Settings.PolicySchema)}) IS NULL "
        + $"EXEC({Literal($"CREATE SCHEMA {Delimit(Settings.PolicySchema)}")})");

    var functionSql =
        $"CREATE FUNCTION {function}(@TenantId uniqueidentifier) "
        + "RETURNS TABLE WITH SCHEMABINDING AS RETURN SELECT 1 AS [Allowed] "
        + $"WHERE @TenantId = CAST(SESSION_CONTEXT({Literal(Settings.SessionKey)}) "
        + "AS uniqueidentifier)";
    Statement(builder, $"IF OBJECT_ID({Literal(function)}, N'IF') IS NULL "
        + $"EXEC({Literal(functionSql)})");

    Statement(builder,
        $"CREATE SECURITY POLICY {PolicyName(table, schema ?? "dbo")} "
        + $"ADD FILTER PREDICATE {predicate} ON {target}, "
        + $"ADD BLOCK PREDICATE {predicate} ON {target} AFTER INSERT, "
        + $"ADD BLOCK PREDICATE {predicate} ON {target} AFTER UPDATE "
        + "WITH (STATE = ON)");
}

protected override void Generate(DropTableOperation operation, IModel? model,
    MigrationCommandListBuilder builder, bool terminate = true)
{
    // SCHEMABINDING: zuerst die Policy, gefunden über die Tabelle, die sie bindet
    DropPolicy(operation.Name, TableSchema(operation.Schema, model), builder);
    base.Generate(operation, model, builder, terminate);
}
\end{sourcecode}

\code{SCHEMABINDING} akzeptiert nur zweiteilige Namen. Eine Tabelle ohne explizites Schema
qualifizieren die Aufrufer deshalb mit dem Default-Schema des Modells oder, falls keines
gesetzt ist, mit \code{dbo}; sonst lehnt SQL Server die Security Policy ab.

\begin{important}[Verschachteltes SQL]
  Das \code{EXEC(\allowbreak{}N'.\allowbreak{}.\allowbreak{}.\allowbreak{}')} enthält SQL in
  einem String-Literal, das wiederum Literale enthält. Die Verschachtelung entsteht
  ausschließlich über \code{Generate\allowbreak{}Sql\allowbreak{}Literal} des
  String-Mappings, das Anführungszeichen korrekt verdoppelt; Bezeichner laufen über
  \code{Delimit\allowbreak{}Identifier}. Session-Schlüssel und Schema werden zusätzlich beim
  Konfigurieren validiert.
\end{important}

\section{Schemaänderungen an geschützten Tabellen}\label{sec:sqlserver-schema-bound-changes}

\code{SCHEMABINDING} schützt die Predicate-Funktion, bindet aber auch die Tabelle. SQL Server
lehnt \code{sp\_\allowbreak{}rename} für Tabelle und Mandantenspalte, \code{ALTER SCHEMA
.\allowbreak{}.\allowbreak{}.\allowbreak{} TRANSFER}, \code{ALTER COLUMN} und \code{DROP COLUMN}
ab, solange eine Security Policy die Tabelle bindet: „object participates in enforced
dependencies“. Pro Operation lässt sich das nicht lösen. EF sortiert die Operationen einer
Migration selbst, und beim Umbenennen der Mandantenspalte kommt die
\code{Rename\allowbreak{}Column\allowbreak{}Operation} \emph{vor} der
\code{Alter\allowbreak{}Table\allowbreak{}Operation}, die die geänderte Annotation trägt. Eine
reine Tabellenumbenennung erzeugt gar keine \code{Alter\allowbreak{}Table\allowbreak{}Operation}.

Der Generator überschreibt deshalb den Einstieg, der die ganze Operationsliste bekommt. Er
sucht alle Tabellen, deren gebundene Objekte sich ändern, verfolgt ihre Umbenennungen und
umklammert die Liste. Vorne entfernt er die Policies der betroffenen Tabellen unter ihren alten
Namen. Hinten legt er sie mit den endgültigen Tabellen-, Schema- und Spaltennamen aus dem
Zielmodell neu an:

\begin{sourcecode}[Generiertes SQL beim Umbenennen der Mandantenspalte (\allowbreak{}gekürzt)]{sql}
DECLARE @policy nvarchar(400) = (SELECT TOP (1) ... FROM sys.security_policies AS p
    WHERE EXISTS (SELECT 1 FROM sys.security_predicates AS s WHERE s.object_id = p.object_id
        AND s.target_object_id = OBJECT_ID(N'[dbo].[Invoice]')
        AND s.predicate_definition LIKE N'%[[]fn[_]tenant[_]predicate]%'));
IF @policy IS NOT NULL EXEC(N'DROP SECURITY POLICY ' + @policy);
EXEC sp_rename N'[Invoice].[TenantId]', N'Mandant', 'COLUMN';
EXEC sp_rename N'[Invoice].[IX_Invoice_TenantId]', N'IX_Invoice_Mandant', 'INDEX';
IF SCHEMA_ID(N'acme') IS NULL EXEC(N'CREATE SCHEMA [acme]');
IF OBJECT_ID(N'[acme].[fn_tenant_predicate]', N'IF') IS NULL EXEC(N'CREATE FUNCTION ...');
CREATE SECURITY POLICY [acme].[tenant_dbo_Invoice]
    ADD FILTER PREDICATE [acme].[fn_tenant_predicate]([Mandant]) ON [dbo].[Invoice], ...
\end{sourcecode}

Drei Entscheidungen stecken darin:

\begin{itemize}
  \item \textbf{Transaktion statt Ausschalten.} EF führt die Befehle einer Migration in einer
    Transaktion aus, und DDL ist in SQL Server transaktional. Scheitert ein Schritt, rollt
    SQL Server auch das Entfernen der Policy zurück. Eine dauerhaft entfernte oder mit
    \code{STATE = OFF} abgeschaltete Policy kann so nicht entstehen. Der Down-Weg ist dieselbe
    Umklammerung mit vertauschten Namen.
  \item \textbf{Entfernen über die Bindung, nicht über den Namen.} Gelöscht wird nur eine
    Policy, deren Prädikat die Predicate-Funktion des Plugins verwendet und die genau diese
    Tabelle bindet. Fremde Policies und Tabellen ohne RLS bleiben unberührt, und Policies
    unter einem früheren Namen werden trotzdem gefunden.
  \item \textbf{Schemaqualifizierte Namen.} Bis 1.0 hieß die Policy
    \code{tenant\_\allowbreak{}Tabelle}, sodass gleichnamige Tabellen in zwei Schemas
    kollidierten. Neue Policies heißen \code{tenant\_\allowbreak{}Schema\_\allowbreak{}Tabelle}.
    Bestehende Installationen behalten den alten Namen, bis die Tabelle das nächste Mal
    angefasst wird; dann ersetzt der Generator die Policy. Die verbleibende Mehrdeutigkeit
    (Schema \code{a\_b} mit Tabelle \code{c} gegen Schema \code{a} mit Tabelle
    \code{b\_c}) und Namen über 128 Zeichen lehnt die Convention beim Modellbau ab.
\end{itemize}

Grenzen: Das Wiederherstellen braucht das Zielmodell. Das ist der Normalfall, weil
\code{dotnet ef migrations add} zu jeder Migration eine Designer-Datei erzeugt. Fehlt es,
oder kommt eine geänderte Tabelle darin nicht vor, bricht der Generator ab, statt eine Policy
fallen zu lassen. Eigene schemagebundene Objekte der Anwendung (Views mit
\code{SCHEMABINDING}, eigene Security Policies) behandelt das Plugin nicht; sie blockieren
dieselben Operationen weiterhin. Datenoperationen, die EF zwischen die Schemaänderungen
sortiert, laufen innerhalb der Transaktion ohne Policy, also mit den Rechten der
Migrationsrolle. \code{Sql\allowbreak{}Server\allowbreak{}Migration\allowbreak{}Tests} prüft
mit von EF erzeugten Diffs gegen einen echten SQL Server: Anlegen, Spalte und Tabelle
umbenennen, Schemawechsel, jeden Schritt zurück und Löschen, nach jedem Schritt Lese- und
Schreibisolation über den App-Login. Dazu kommen gleichnamige Tabellen in zwei Schemas, eine
Policy unter dem alten Namen und die fremde Policy einer Tabelle ohne RLS.

\section{Session-Schlüssel und Policy-Schema ändern}\label{sec:sqlserver-session-settings}

\code{session\allowbreak{}Key} und \code{policy\allowbreak{}Schema} sind Options, wirken aber
an zwei Stellen: zur Laufzeit im Session-Interceptor, der den Schlüssel setzt, und in der
Datenbank in der Predicate-Funktion, die ihn liest. Bis 1.0 las der Generator beide Werte aus
den \emph{aktuellen} Options und legte die Funktion nur an, wenn sie fehlte. Ein Wechsel von
\code{TenantId} auf \code{NewTenant} erzeugte deshalb keine Migration. Der Interceptor setzte
danach den neuen Schlüssel, die Funktion las weiter den alten, und kein Mandant sah mehr Zeilen.
Schlimmer noch: Eine alte Migration hätte beim erneuten Ausführen, etwa für eine neue
Testdatenbank, Objekte mit den neuen Werten erzeugt.

Seit 1.1 trennt das Paket die beiden Rollen. Die Convention schreibt Schlüssel und Schema als
Modell-Annotationen \code{Acme:\allowbreak{}Sql\allowbreak{}Server\allowbreak{}Rls:\allowbreak{}Session\allowbreak{}Key}
und \code{Acme:\allowbreak{}Sql\allowbreak{}Server\allowbreak{}Rls:\allowbreak{}Policy\allowbreak{}Schema}
ins Modell und damit in den Snapshot. Der Annotation-Provider gibt sie an das relationale Modell
weiter, und der Differ erzeugt bei einer Änderung eine
\code{Alter\allowbreak{}Database\allowbreak{}Operation} mit altem und neuem Wert. Der Generator
liest beide Werte aus dem Zielmodell der Migration, die er gerade erzeugt; die Options gelten nur
noch für Modelle, die die Werte nicht festhalten. Damit bleibt jede Migration bei dem Stand,
mit dem sie erzeugt wurde.

\begin{datatable}{H{0.70}H{1.30}}
  \textbf{Fall} & \textbf{Verhalten der Migration} \\
  \tablehead
  Unverändert & Keine Operation, keine Phantomänderung \\
  \tablerule
  Schlüssel oder Schema geändert, alter Wert im Snapshot & In einer Transaktion: Policies aller
  geschützten Tabellen entfernen, alte Funktion löschen, Funktion mit den neuen Werten anlegen,
  Policies neu anlegen. Der Down-Weg macht dasselbe rückwärts. \\
  \tablerule
  Erste Migration eines Contexts oder erste nach dem Upgrade von 1.0 (alter Wert unbekannt) &
  Nichts wird gelöscht. Eine vorhandene Funktion muss den Schlüssel lesen, den das Modell nennt,
  sonst bricht die Migration mit \code{THROW 50000} ab. \\
  \tablerule
  Down auf einen Snapshot vor 1.1 & Funktion und Policies bleiben, wie sie sind. \\
\end{datatable}

Die dritte Zeile ist eine bewusste Grenze. Ohne festgehaltenen alten Wert lässt sich nicht
beweisen, dass eine vorhandene Funktion dem Plugin gehört und nicht einem anderen Context.
Eine bestehende Installation wechselt den Schlüssel deshalb in zwei Schritten: erst die erste
Migration nach dem Upgrade mit unverändertem Schlüssel, dann der Wechsel in einer eigenen
Migration.

Dieselbe Prüfung schützt mehrere Contexts in einer Datenbank. Die Funktion ist pro
Policy-Schema eine, und jede Policy des Plugins prüft nach dem Anlegen der Funktion deren
Definition. Konfiguriert ein zweiter Context einen anderen Schlüssel im selben Schema, scheitert
seine Migration, statt seine Tabellen an eine Funktion zu binden, die den falschen Schlüssel
liest. Contexts mit verschiedenen Schlüsseln brauchen verschiedene Policy-Schemas.

Laufzeit und Datenbank können trotzdem kurz auseinanderlaufen. Läuft eine neue App-Version mit
neuem Schlüssel, bevor die Migration eingespielt ist, setzt der Interceptor einen Schlüssel, den
die Funktion nicht liest. Die Folge ist kein Datenleck, sondern \emph{fail closed}: keine Zeilen,
und Schreibversuche scheitern am BLOCK-Prädikat. Die Migration gehört deshalb vor die neue
App-Version. \code{Sql\allowbreak{}Server\allowbreak{}Session\allowbreak{}Settings\allowbreak{}Tests}
prüft ohne Datenbank die Annotationen, den Snapshot, den leeren Diff, die erkannte Änderung und
die historische Migration. Gegen einen echten SQL Server prüft es den Wechsel von Schlüssel und
Schema mit Up und Down, das Upgrade einer 1.0-Installation und den zweiten Context.

\section{Laufzeit: derselbe Ablauf, anderer Befehl}

% source: src/Acme.EntityFrameworkCore.Tenancy.SqlServer/Runtime/SqlServerTenantSessionInterceptor.cs
\begin{sourcecode}[Runtime/\allowbreak{}Sql\allowbreak{}Server\allowbreak{}Tenant\allowbreak{}Session\allowbreak{}Interceptor.\allowbreak{}cs]{csharp}
public class SqlServerTenantSessionInterceptor : TenantSessionInterceptor
{
    public SqlServerTenantSessionInterceptor(ITenantProvider tenants)
        : base(tenants)
    {
    }

    protected override DbCommand CreateSessionCommand(
        DbConnection connection, DbContext? context, TenantId? tenant)
    {
        var key = context?.GetService<ISqlServerTenancySingletonOptions>()
                      .Settings.SessionKey
                  ?? SqlServerTenancySettings.Default.SessionKey;

        var command = connection.CreateCommand();
        command.CommandText = "EXEC sp_set_session_context @key = @key, @value = @value";
        AddParameter(command, "@key", key);
        AddParameter(command, "@value", tenant?.Value);    // NULL: keine Zeilen sichtbar
        return command;
    }
}
\end{sourcecode}

\section{Bewusst nicht portiert}

Das SQL-Server-Paket ist ein eigenständiger, aber schlankerer Provider. Fluent-API-Policies,
der optionale Query-Filter, eigene SQL-Knoten, Aggregate, Snapshot-Code-Generator und
Scaffolding fehlen dort. Das ist eine Produktentscheidung, kein technisches Hindernis: Jedes
dieser Kapitel lässt sich mit denselben Mustern auf SQL Server übertragen. Die Liste gehört
in die Dokumentation des Pakets, damit Nutzer wissen, was sie erwarten können.

% source: tests/Acme.EntityFrameworkCore.Tenancy.SqlServer.Tests/SqlServerTests.cs
\begin{sourcecode}[tests/\allowbreak{}.\allowbreak{}.\allowbreak{}.\allowbreak{}Sql\allowbreak{}Server.\allowbreak{}Tests/\allowbreak{}Sql\allowbreak{}Server\allowbreak{}Tests.\allowbreak{}cs (\allowbreak{}Auszug)]{csharp}
[Fact]
public async Task Block_predicate_rejects_rows_of_other_tenant()
{
    await using var db = ContextFor(TenantId.New());

    var ex = await Assert.ThrowsAsync<SqlException>(() => db.Database.ExecuteSqlAsync(
        $"INSERT INTO [Orders] ([TenantId], [Total]) VALUES ({Guid.NewGuid()}, 1)"));

    Assert.Equal(33504, ex.Number);      // Verstoß gegen BLOCK PREDICATE
}
\end{sourcecode}

Die Tests laufen wie bei PostgreSQL in zwei Ebenen: ohne Datenbank für Modell, Abfrage- und
Migrations-SQL und mit Testcontainers gegen einen echten SQL Server. Dort verbindet sich die
Anwendung mit einem eigenen Login statt als \code{sa}, um die Produktionssituation
nachzustellen.
